\documentclass[a4paper,11pt]{article}

%%% Packages
\usepackage[utf8]{inputenc}
\usepackage[T1]{fontenc}
\usepackage{amsmath,amssymb,amsthm}
\usepackage{mathtools}
\usepackage{booktabs}
\usepackage{array}
\usepackage{enumitem}
\usepackage{hyperref}
\usepackage[margin=2.5cm]{geometry}
\usepackage{xcolor}

\hypersetup{colorlinks=true,linkcolor=blue!45!black,citecolor=blue!45!black,urlcolor=blue!45!black}

%%% Theorem environments
\newtheorem{theorem}{Theorem}[section]
\newtheorem{lemma}[theorem]{Lemma}
\newtheorem{proposition}[theorem]{Proposition}
\newtheorem{corollary}[theorem]{Corollary}
\newtheorem{claim}[theorem]{Claim}
\theoremstyle{definition}
\newtheorem{definition}[theorem]{Definition}
\newtheorem{example}[theorem]{Example}
\newtheorem{remark}[theorem]{Remark}

%%% Notation macros
\newcommand{\ALCI}{\ensuremath{\mathcal{ALCI}}}
\newcommand{\ALCIRCC}[1]{\ensuremath{\mathcal{ALCI}_{\mathit{RCC#1}}}}
\newcommand{\DR}{\ensuremath{\mathsf{DR}}}
\newcommand{\PO}{\ensuremath{\mathsf{PO}}}
\newcommand{\EQ}{\ensuremath{\mathsf{EQ}}}
\newcommand{\PP}{\ensuremath{\mathsf{PP}}}
\newcommand{\PPI}{\ensuremath{\mathsf{PPI}}}
\newcommand{\inv}{\ensuremath{\mathsf{inv}}}
\newcommand{\comp}{\ensuremath{\mathsf{comp}}}
\newcommand{\cl}{\ensuremath{\mathsf{cl}}}
\newcommand{\tp}{\ensuremath{\mathsf{tp}}}
\newcommand{\Tp}{\ensuremath{\mathsf{Tp}}}
\newcommand{\EXPTIME}{\textsc{ExpTime}}
\newcommand{\Desc}{\ensuremath{\mathsf{Desc}}}
\newcommand{\Core}{\ensuremath{\mathsf{Core}}}
\newcommand{\Union}{\ensuremath{\mathsf{Union}}}
\newcommand{\Safe}{\ensuremath{\mathsf{Safe}}}
\newcommand{\interp}{\ensuremath{\mathcal{I}}}
\newcommand{\DeltaI}{\ensuremath{\Delta^{\mathcal{I}}}}

\setlength{\parindent}{0pt}
\setlength{\parskip}{0.55em}


\title{\textbf{Decidability of $\ALCIRCC{5}$ via\\Two-Tier Quotient Construction}\\[0.4em]
\large Period descriptors, PP-kernel nodes,\\
and the full tractability of RCC5\\[0.3em]
\normalsize (Fourth revision, April 2026)}
\author{Claude (Anthropic)\\
\small Prompted by Michael Wessel}
\date{April 2026}

\begin{document}
\maketitle

\begin{center}
\fbox{\parbox{0.92\textwidth}{%
\textbf{Disclaimer.}\; This document was generated by Claude
(Anthropic), prompted by Michael Wessel.  The results have not been
independently verified.  We invite scrutiny and corrections.

\smallskip\textbf{Revision history.}\; This is the fourth
revision.  The first version was reviewed by GPT-5.4
Pro~\cite{GPTReview2026}, who raised five objections; these were
addressed in the second version~\cite{ClaudeResponse2026}.  GPT-5.4's
second review~\cite{GPTReview2026b} raised four further objections,
addressed in the third version.  GPT-5.4's third
review~\cite{GPTReview2026c} raised four more objections.  This
fourth version incorporates all fixes: a new phasewise safety
validity condition ($V_{\mathrm{safe}}$); exact-relation extraction
proofs via backward forcing ($\DR$, $\PP$) and forward absorption
($\PPI$); a constructive quotient approach that breaks the
size-bound circularity; and an honest acknowledgment of a remaining
$\PO$ gap.}}
\end{center}

\begin{abstract}
We study the decidability of concept satisfiability in
$\ALCIRCC{5}$---the description logic $\ALCI$ extended with an RCC5
role box.  We introduce a finite \emph{two-tier quotient}: Tier~1
captures each infinite $\PP$-chain's eventually periodic tail as a
\emph{period descriptor}, while Tier~2 captures cross-chain
interactions through \emph{PP-kernel nodes}.  A valid quotient can
be unfolded into a genuine model using the \emph{full tractability}
of RCC5 (Renz \& Nebel 1999).

For demands involving $\DR$, $\PP$, and $\PPI$, we prove complete
decidability.  The key ingredients are: (i)~a \emph{backward
forcing} property for $\DR$ and $\PP$ (the composition table forces
these relations backward along $\PP$-chains, guaranteeing all-phase
type-safety of witness types); (ii)~\emph{forward absorption} for
$\PPI$ ($\comp(\PPI, \PPI) = \{\PPI\}$ ensures stabilized witnesses);
(iii)~a \emph{constructive quotient} with one kernel per descriptor,
giving a computable size bound with no circularity; and (iv)~a
\emph{chain-unfolding lift lemma} bridging the finite quotient to
the infinite model.

For $\PO$, a genuine gap remains: the composition table provides
neither backward forcing nor forward absorption, and we exhibit a
realizable descriptor where the constant-interface architecture is
provably incomplete.  We precisely characterize the \emph{PO-coherent
fragment} where the proof succeeds and identify the $\PO$ obstacle
as the key remaining open problem.
\end{abstract}


%% ============================================================
\section{Introduction}\label{sec:intro}

The description logic $\ALCIRCC{5}$ extends $\ALCI$ with five jointly
exhaustive and pairwise disjoint (JEPD) base relations from the
Region Connection Calculus RCC5~\cite{RandellCuiCohn1992}: $\DR$,
$\PO$, $\EQ$, $\PP$, $\PPI$.  Wessel~\cite{Wessel2003} introduced
this logic and left its decidability open.  The logic lacks both the
finite model property and the tree model property: the concept
$C_\infty = (\exists\PP.\top) \sqcap (\forall\PP.\exists\PP.\top)$ is
satisfiable only in infinite models containing unbounded $\PP$-chains.

Several recent approaches have been attempted:
a quasimodel method~\cite{Claude2026} (gap in completeness),
a contextual tableau~\cite{GPT2026} (whose completeness conjecture
$\mathsf{FW}(C,N)$ is false~\cite{ClaudeFW2026}), and a
triangle-type approach~\cite{ClaudeTriangle2026} (conditional on an
unproven conjecture).  All share a common difficulty: assigning
globally consistent RCC5 edges in complete-graph models.

\textbf{Results.}\; We prove decidability of concept satisfiability
in $\ALCIRCC{5}$ for the \emph{PO-coherent fragment}
(Definition~\ref{def:po-coherent}), which includes all concepts
without $\forall\PO$ restrictions and all concepts where the
$\forall\PO$ restrictions are uniform across phases.  For the full
logic, we identify a precisely characterized $\PO$ obstacle
(\S\ref{sec:po-gap}) and leave its resolution as the main open
problem.


%% ============================================================
\section{Preliminaries}\label{sec:prelim}

\subsection{The logic $\ALCIRCC{5}$}

Concepts are built from concept names, the RCC5 base relations
$R \in \{\DR, \PO, \EQ, \PP, \PPI\}$ as role names, and the
constructors $\sqcap$, $\sqcup$, $\neg$, $\exists R.C$,
$\forall R.C$ in negation normal form (NNF).  We adopt
\textbf{strong $\EQ$ semantics}.

An interpretation $\interp = (\DeltaI, \cdot^{\interp})$ assigns a
total function $\rho^{\interp}: \DeltaI \times \DeltaI \to
\{\DR, \PO, \EQ, \PP, \PPI\}$ subject to: (i)~$\rho(d,d) = \EQ$;
(ii)~$\rho(d,e) = \inv(\rho(e,d))$; (iii)~composition consistency:
$\rho(d_1, d_3) \in \comp(\rho(d_1, d_2), \rho(d_2, d_3))$ for all
$d_1, d_2, d_3$.

\subsection{Types, closure, composition table}

Given $C_0$ in NNF, $\cl(C_0)$ is the closure under subformulas.
A \emph{Hintikka type} $\tau \subseteq \cl(C_0)$ is a maximal
propositionally consistent subset.  Write $\Tp(C_0)$ for the set of
types; $|\Tp(C_0)| \le 2^{|\cl(C_0)|}$.

The RCC5 composition table
$\comp: \{\DR,\PO,\PP,\PPI\}^2 \to 2^{\{\DR,\PO,\PP,\PPI\}}$ (with
$\EQ$ excluded for distinct elements) is:

\begin{center}
\renewcommand{\arraystretch}{1.15}
\begin{tabular}{c|cccc}
$\comp(R_1, R_2)$ & $\DR$ & $\PO$ & $\PP$ & $\PPI$ \\
\hline
$\DR$ & \{all\} & \{\DR,\PO,\PP\} & \{\DR,\PO,\PP\} & \{\DR\} \\
$\PO$ & \{\DR,\PO,\PPI\} & \{all\} & \{\PO,\PP\} & \{\DR,\PO,\PPI\} \\
$\PP$ & \{\DR\} & \{\DR,\PO,\PP\} & \{\PP\} & \{all\} \\
$\PPI$ & \{\DR,\PO,\PPI\} & \{\PO,\PPI\} & \{\PO,\PP,\PPI\} & \{\PPI\} \\
\end{tabular}
\end{center}

\subsection{Type-safety for RCC5 relations}\label{sec:type-safety}

\begin{definition}[Type-safe relations]\label{def:safe}
For Hintikka types $\tau_1, \tau_2 \in \Tp(C_0)$, define
\[
\Safe(\tau_1, \tau_2) = \{R \in \{\DR,\PO,\PP,\PPI\} :
\text{$R$ respects all universals in $\tau_1$ and $\tau_2$}\}
\]
where $R$ \emph{respects the universals} means: for every
$\forall R.C \in \tau_1$, $C \in \tau_2$; and for every
$\forall\inv(R).C \in \tau_2$, $C \in \tau_1$.
\end{definition}

\begin{lemma}[Model relations are type-safe]\label{lem:model-safe}
In any model $\interp$, for distinct $d, e \in \DeltaI$:
$\rho(d, e) \in \Safe(\tp(d), \tp(e))$.
\end{lemma}

\begin{proof}
Immediate from the semantics: if $\forall R.C \in \tp(d)$ and
$\rho(d, e) = R$, then $e \in C^{\interp}$, so $C \in \tp(e)$.
The inverse direction is symmetric.
\end{proof}

\subsection{Full tractability of RCC5}

\begin{theorem}[Full RCC5 tractability, Renz \& Nebel
\protect{\cite{RenzNebel1999}}, Renz
\protect{\cite{Renz1999}}]\label{thm:full-tractability}
Every path-consistent disjunctive RCC5 constraint network over the
base relations $\{\DR, \PO, \EQ, \PP, \PPI\}$ is satisfiable.
Equivalently: if path-consistency enforcement does not empty any
domain, then a consistent atomic refinement exists.\qed
\end{theorem}

This is the ``patchwork property'' of RCC5: local consistency
(path-consistency) implies global consistency.


%% ============================================================
\section{Algebraic properties of $\PP$}\label{sec:pp-algebra}

\begin{theorem}[Universal self-absorption of $\PP$]
\label{thm:self-absorbing}
\begin{enumerate}[label=(\alph*)]
\item \textbf{Left absorption:} $R \in \comp(\PP, R)$ for every
  $R \in \{\DR,\PO,\PP,\PPI\}$.
\item \textbf{Right absorption:} $R \in \comp(R, \PP)$ for every $R$.
\item \textbf{PPI right-absorption:} $R \in \comp(\PPI, R)$ for every $R$.
\item \textbf{PPI left-absorption:} $R \in \comp(R, \PPI)$ for every $R$.
\end{enumerate}
\end{theorem}

\begin{proof}
Direct verification against the composition table.
\end{proof}

\begin{corollary}[Reflexive $\PP$ is composition-consistent]
\label{cor:reflexive-pp}
A node $k$ with $\rho(k,k) = \PP$ satisfies all composition
constraints with any external node $m$.\qed
\end{corollary}

\begin{theorem}[$\PP$-transitivity]
\label{thm:no-pp-cycle}
$\comp(\PP, \PP) = \{\PP\}$.  Consequently, $\PP$ on any finite set
of elements is a strict partial order.\qed
\end{theorem}

\begin{lemma}[Backward forcing]\label{lem:backward}
Let $(d_i)_{i \ge 0}$ be a $\PP$-chain with $\rho(d_j, w) = R$ for
some element $w$ and position $j$.
\begin{enumerate}[label=(\alph*)]
\item If $R = \DR$: then $\rho(d_i, w) = \DR$ for all $i \le j$.
\item If $R = \PP$: then $\rho(d_i, w) = \PP$ for all $i \le j$.
\end{enumerate}
\end{lemma}

\begin{proof}
For $i < j$: $\rho(d_i, d_j) = \PP$, so
$\rho(d_i, w) \in \comp(\PP, \rho(d_j, w)) = \comp(\PP, R)$.
Now $\comp(\PP, \DR) = \{\DR\}$ and $\comp(\PP, \PP) = \{\PP\}$:
both are singletons forcing the same value backward.
\end{proof}

\begin{remark}[Why backward forcing fails for $\PO$ and $\PPI$]
\label{rem:no-backward-po}
$\comp(\PP, \PO) = \{\DR, \PO, \PP\}$ and
$\comp(\PP, \PPI) = \{$all$\}$ are not singletons.  Hence
$\PO$ and $\PPI$ are \emph{not} backward-forced: knowing
$\rho(d_j, w) = \PO$ does not determine $\rho(d_i, w)$ for
$i < j$.  This asymmetry is the algebraic root of the $\PO$ gap
identified in \S\ref{sec:po-gap}.
\end{remark}

\begin{lemma}[Forward transition table]\label{lem:forward}
Along a $\PP$-chain $(d_i)$, the valid transitions for
$\rho(d_i, w) \to \rho(d_{i+1}, w)$ (satisfying \emph{both} the
forward constraint
$\rho(d_{i+1}, w) \in \comp(\PPI, \rho(d_i, w))$ and the backward
constraint
$\rho(d_i, w) \in \comp(\PP, \rho(d_{i+1}, w))$) are:
\begin{center}
\renewcommand{\arraystretch}{1.15}
\begin{tabular}{cl}
\toprule
\textbf{Current} & \textbf{Valid next values} \\
\midrule
$\DR$ & $\DR$, $\PO$, $\PPI$ \\
$\PO$ & $\PO$, $\PPI$ \\
$\PP$ & $\PO$, $\PP$, $\PPI$ \\
$\PPI$ & $\PPI$ \\
\bottomrule
\end{tabular}
\end{center}
In particular: (i)~$\PPI$ is forward-absorbing; (ii)~once the
sequence leaves $\DR$ (to $\PO$ or $\PPI$), it never returns to
$\DR$; (iii)~once it leaves $\PO$ (to $\PPI$), it never returns
to $\PO$.
\end{lemma}

\begin{proof}
The forward constraint gives:
$\comp(\PPI, \DR) = \{\DR,\PO,\PPI\}$,
$\comp(\PPI, \PO) = \{\PO,\PPI\}$,
$\comp(\PPI, \PP) = \{\PO,\PP,\PPI\}$,
$\comp(\PPI, \PPI) = \{\PPI\}$.
The backward constraint eliminates further values.  For instance,
from $\DR$: the forward constraint allows $\DR, \PO, \PPI$; checking
$\DR \in \comp(\PP, R')$ for each candidate $R'$: $\comp(\PP, \DR) =
\{\DR\}$ ($\DR \in$~\checkmark), $\comp(\PP, \PO) = \{\DR, \PO, \PP\}$
($\DR \in$~\checkmark), $\comp(\PP, \PPI) = \{$all$\}$
($\DR \in$~\checkmark).  All survive.  From $\PO$: forward allows
$\PO, \PPI$; backward check: $\PO \in \comp(\PP, \PO) =
\{\DR, \PO, \PP\}$ (\checkmark), $\PO \in \comp(\PP, \PPI) =
\{$all$\}$ (\checkmark).  Both survive.  The no-return properties
follow by inspection.
\end{proof}


%% ============================================================
\section{$\PP$-chain structure}\label{sec:chain}

\begin{definition}[$\PP$-chain]
In a model $\interp$, a \emph{$\PP$-chain} is a sequence
$(d_i)_{i \ge 0}$ with $\rho(d_i, d_{i+1}) = \PP$ for all $i$.
By Theorem~\ref{thm:no-pp-cycle}, $\rho(d_i, d_j) = \PP$ for $i < j$.
\end{definition}

\begin{lemma}[External relation stabilization]\label{lem:external}
For any element $e$ and chain $(d_i)$, the sequence
$\rho(d_i, e)$ eventually stabilizes to a constant value.
\end{lemma}

\begin{proof}
By Lemma~\ref{lem:forward}, the sequence takes values in
$\{\DR, \PO, \PP, \PPI\}$ with only the listed transitions.
The non-return properties (once leaving $\DR$, $\PO$, or $\PP$,
the sequence cannot return) ensure eventual constancy.
\end{proof}

\begin{definition}[Recurrent types, $T_\infty$]\label{def:tinf}
For a $\PP$-chain $(d_i)_{i \ge 0}$, define
$T_\infty = \{\tau \in \Tp(C_0) : \tau = \tp(d_i) \text{ for
infinitely many } i\}$.
\end{definition}

\begin{lemma}[$T_\infty$ is non-empty and eventually
exhaustive]\label{lem:tinf}
$T_\infty \ne \emptyset$.  There exists $M$ such that
$\tp(d_i) \in T_\infty$ for all $i \ge M$.
\end{lemma}

\begin{proof}
Since $|\Tp(C_0)|$ is finite and the sequence $(\tp(d_i))_{i \ge 0}$
is infinite, some type appears infinitely often (pigeonhole), so
$T_\infty \ne \emptyset$.  Types not in $T_\infty$ appear only
finitely often; let $M$ be the last index at which any such type
appears.
\end{proof}

\begin{lemma}[Phase-uniform type-safety]\label{lem:phase-uniform}
Let $(d_i)$ be a $\PP$-chain in a model $\interp \models C_0$
with $T_\infty = \{\tau_1, \ldots, \tau_p\}$.  Let $w$ be any
element with stabilized relation $S = \rho(d_i, w)$ for all
sufficiently large $i$.  Then
$S \in \Safe(\tau_j, \tp(w))$ for every $\tau_j \in T_\infty$.
\end{lemma}

\begin{proof}
Fix $\tau_j \in T_\infty$.  Choose $k$ large enough that both
$\tp(d_k) = \tau_j$ and $\rho(d_k, w) = S$.  For any
$\forall S.C \in \tau_j$: the model satisfies $\forall S.C$ at $d_k$,
and $\rho(d_k, w) = S$, so $C \in \tp(w)$.  The inverse direction
is symmetric.  Hence $S \in \Safe(\tau_j, \tp(w))$.
\end{proof}

\begin{lemma}[Backward forcing gives all-phase safety]
\label{lem:backward-safe}
Let $(d_i)$ be a $\PP$-chain, $T_\infty = \{\tau_1, \ldots, \tau_p\}$,
and $R \in \{\DR, \PP\}$.  If $\rho(d_j, w) = R$ for some position
$j > M$ (past the exhaustion index), then
$R \in \Safe(\tau_k, \tp(w))$ for every $\tau_k \in T_\infty$.
\end{lemma}

\begin{proof}
By Lemma~\ref{lem:backward}, $\rho(d_i, w) = R$ for all $i \le j$.
Since $j > M$, every $\tau_k \in T_\infty$ appears at some position
$i \le j$.  At that position, $\rho(d_i, w) = R$ and
$\tp(d_i) = \tau_k$, so $R \in \Safe(\tau_k, \tp(w))$ by
Lemma~\ref{lem:model-safe}.
\end{proof}


%% ============================================================
\section{Period descriptors}\label{sec:period}

\begin{definition}[Period descriptor]\label{def:period-descriptor}
A \emph{period descriptor} for $C_0$ is a tuple
$\Desc = (\tau_1, \ldots, \tau_p)$ with $p \ge 1$ and each
$\tau_j \in \Tp(C_0)$, satisfying:
\begin{enumerate}[label=(V\arabic*)]
\item Each $\tau_j$ is a Hintikka type in $\cl(C_0)$.
\item \textbf{$\forall\PP$-closure:} $\forall\PP.C \in \tau_i$
  for some $i$ implies $C \in \tau_j$ for all $j$.
\item \textbf{$\forall\PPI$-closure:} $\forall\PPI.C \in \tau_i$
  for some $i$ implies $C \in \tau_j$ for all $j$.
\item \textbf{Chain continuation:} For each $i$, there exists
  $\exists\PP.C \in \tau_i$ such that $C \in \tau_j$ for some~$j$.
\end{enumerate}
\end{definition}

\begin{theorem}[$T_\infty$ yields a valid period descriptor]
\label{thm:tinf-descriptor}
Let $(d_i)_{i \ge 0}$ be a $\PP$-chain in a model $\interp \models C_0$.
Any enumeration of $T_\infty = \{\tau_1, \ldots, \tau_p\}$ is a valid
period descriptor.\qed
\end{theorem}

\begin{definition}[Stabilized core and union]\label{def:core-union}
For $\Desc = (\tau_1, \ldots, \tau_p)$:
\begin{itemize}
\item $\Core(\Desc) = \bigcap_{j=1}^{p} \tau_j$.
\item $\Union(\Desc) = \bigcup_{j=1}^{p} \tau_j$.
\end{itemize}
\end{definition}

\begin{definition}[PO-coherent descriptor]\label{def:po-coherent}
A period descriptor $\Desc = (\tau_1, \ldots, \tau_p)$ is
\emph{PO-coherent} if for every
$\exists\PO.C \in \Union(\Desc)$, there exists a Hintikka type
$\sigma$ with $C \in \sigma$ and
$\PO \in \bigcap_{j=1}^p \Safe(\tau_j, \sigma)$.

A concept $C_0$ is \emph{PO-coherent} if every realizable period
descriptor (i.e., one extractable as $T_\infty$ from a model of $C_0$)
is PO-coherent.
\end{definition}

\begin{remark}[When PO-coherence holds automatically]
\label{rem:po-coherent}
A descriptor is automatically PO-coherent if: (a)~it contains no
$\exists\PO$ demands, or (b)~for every $\forall\PO.D$ appearing in
\emph{any} phase, $D \in \Core(\Desc)$ (so the $\forall\PO$
universals are the same across all phases), or (c)~no phase contains
any $\forall\PO$ restriction at all.  The PO-incoherent case requires
$\exists\PO.A$ in one phase and $\forall\PO.\neg A$ in a different
phase---a relatively specific configuration.
\end{remark}


%% ============================================================
\section{The two-tier quotient}\label{sec:quotient}

\begin{definition}[Two-tier quotient]\label{def:quotient}
A \emph{two-tier quotient} $\mathcal{Q}$ for $C_0$ consists of:

\medskip\noindent\textbf{Tier~1 (within-chain structure):}
A finite set of PO-coherent period descriptors
$\{\Desc_1, \ldots, \Desc_K\}$, each satisfying (V1)--(V4).

\medskip\noindent\textbf{Tier~2 (cross-chain structure):}
\begin{enumerate}[label=(T\arabic*)]
\item A set of \emph{PP-kernel nodes}
  $\{k_1, \ldots, k_K\}$, one per descriptor.
  Each $k_\alpha$ is associated with descriptor $\Desc_\alpha$.
  We set $\tp(k_\alpha) = \Core(\Desc_\alpha)$ and
  $\rho(k_\alpha, k_\alpha) = \PP$ (bookkeeping device; see
  Remark~\ref{rem:certificate}).
\item A set of \emph{regular nodes} $\{m_1, \ldots, m_L\}$ with
  types $\tp(m_\beta) \in \Tp(C_0)$.
\item An atomic RCC5 labeling $\rho$ on all pairs of
  kernel and regular nodes, satisfying composition consistency (V5)
  and phasewise type-safety ($V_{\mathrm{safe}}$).
\item \textbf{Designated witnesses for kernel demands.}\;
  For each kernel node $k_\alpha$ and each non-$\PP$ existential
  demand $\exists R.C \in \Union(\Desc_\alpha)$ with
  $R \ne \PP$: a designated witness $w$ among the kernel/regular
  nodes, such that $C \in \tp(w)$,
  $\rho(k_\alpha, w) = R$, and
  \begin{equation}\label{eq:t4-allphase}
  R \in \Safe(\tau_j, \tp(w)) \quad\text{for \emph{every} phase
  type $\tau_j$ of $\Desc_\alpha$}.
  \end{equation}
\item For each regular node $m_\beta$ and each existential demand
  $\exists R.C \in \tp(m_\beta)$: a designated witness.
\end{enumerate}

\medskip\noindent\textbf{Validity conditions:}
\begin{enumerate}[label=(V\arabic*),start=5]
\setcounter{enumi}{4}
\item \textbf{Composition consistency:} The atomic network on
  kernel and regular nodes is composition-consistent (CC).
\item \textbf{Off-chain extension:} For each
  $\exists\PP.C \in \Union(\Desc_\alpha)$
  not satisfied within $\Desc_\alpha$ (i.e., $C \notin \tau_j$ for
  any $j$), there is a witness node $w$ with
  $\rho(k_\alpha, w) = \PP$, $C \in \tp(w)$, and
  $\PP \in \Safe(\tau_j, \tp(w))$ for every phase type $\tau_j$.
  The disjunctive constraint network obtained by adding all such
  witnesses simultaneously (with domains
  \begin{equation}\label{eq:v6-domain}
  D(w, e) = \comp(\text{demanded}, \rho(\text{parent}, e)) \;\cap\;
  \Safe(\tp(w), \tp(e))
  \end{equation}
  for all existing nodes $e$) is path-consistent after
  arc-consistency enforcement.
\end{enumerate}

\medskip\noindent\textbf{Phasewise type-safety:}
\begin{enumerate}[label=($V_{\mathrm{safe}}$)]
\item For every kernel $k_\alpha$ and every
  external node $e$ (kernel or regular):
  \[
  \rho(k_\alpha, e) \;\in\;
    \bigcap_{\tau \in \Desc_\alpha}
    \Safe(\tau, \tp(e)).
  \]
  For kernel--kernel edges, the condition applies on both sides.
\end{enumerate}
\end{definition}

\begin{remark}[The quotient is a certificate, not a model]
\label{rem:certificate}
The two-tier quotient is \emph{not} itself an RCC5 network in the
strong-$\EQ$ sense: kernel nodes have reflexive $\PP$-loops.
The quotient is an abstract certificate from which a genuine
strong-$\EQ$ model can be constructed
(Theorem~\ref{thm:soundness}).  During chain unfolding, each
kernel is replaced by infinitely many distinct elements with $\EQ$
self-loops.
\end{remark}

\begin{remark}[All-phase type-safety checks]
\label{rem:allphase}
The type-safety conditions in T4~\eqref{eq:t4-allphase}, V6,
and $V_{\mathrm{safe}}$ are checked against \emph{every} phase type
$\tau_j$ of $\Desc_\alpha$.  This is necessary because the
soundness proof assigns the kernel's relation uniformly to all chain
positions.  Lemma~\ref{lem:phase-uniform} guarantees that
model-derived relations always pass this all-phase check.
\end{remark}


%% ============================================================
\section{Completeness: model implies quotient}
\label{sec:completeness}

\begin{theorem}[Constructive quotient extraction]
\label{thm:completeness}
If $C_0$ is satisfiable and PO-coherent, there exists a valid
two-tier quotient $\mathcal{Q}$ for $C_0$ of size at most
$f(|C_0|)$, where $f$ is a computable function.
\end{theorem}

\begin{proof}
Let $\interp \models C_0$.  Form the tree unraveling $\interp^T$.

\medskip\noindent\textit{Step 1: One kernel per descriptor.}\;
For each maximal $\PP$-chain in $\interp^T$, compute $T_\infty$.
By Theorem~\ref{thm:tinf-descriptor}, $T_\infty$ is a valid period
descriptor.  Group chains by descriptor.  Select one representative
chain per distinct descriptor.  This gives $K \le D$ kernels, where
$D \le 2^{|\Tp(C_0)|} \le 2^{2^{|\cl(C_0)|}}$ is a computable
bound on the number of distinct descriptors.

\medskip\noindent\textit{Step 2: Kernel-to-kernel relations.}\;
For each pair of representative chains $\alpha, \beta$: the doubly
stabilized relation $\rho(k_\alpha, k_\beta)$ is
$\lim_{i,j \to \infty} \rho(d_i^{(\alpha)}, d_j^{(\beta)})$.
This is well-defined by Lemma~\ref{lem:external} applied twice.

By Lemma~\ref{lem:phase-uniform},
$\rho(k_\alpha, k_\beta) \in \Safe(\tau, \Core(\Desc_\beta))$ for
every $\tau \in \Desc_\alpha$ (and symmetrically).  So
$V_{\mathrm{safe}}$ holds for kernel--kernel edges.  Composition
consistency follows from the model (choose sufficiently large
indices).

\medskip\noindent\textit{Step 3: Designated witnesses.}\;
For each kernel $k_\alpha$ and each
$\exists R.C \in \Union(\Desc_\alpha)$ with $R \ne \PP$:

\textbf{Case $R = \DR$:}\; Choose a phase $\tau_j$ containing
$\exists\DR.C$ and a position $i_n > M$ with $\tp(d_{i_n}) = \tau_j$.
The model provides a $\DR$-witness $w$ with $\rho(d_{i_n}, w) = \DR$
and $C \in \tp(w)$.  By Lemma~\ref{lem:backward-safe},
$\DR \in \Safe(\tau_k, \tp(w))$ for all $\tau_k \in T_\infty$.
Create a quotient witness of type $\tp(w)$ with
$\rho(k_\alpha, w') = \DR$.  The all-phase check~\eqref{eq:t4-allphase}
passes.

\textbf{Case $R = \PPI$:}\; At a position $i_n$ past the
stabilization threshold for some witness $w$ with
$\rho(d_{i_n}, w) = \PPI$ and $C \in \tp(w)$: the forward
absorption of $\PPI$ (Lemma~\ref{lem:forward}) gives $S_w = \PPI$.
By Lemma~\ref{lem:phase-uniform}, all-phase safety holds.

\textbf{Case $R = \PO$:}\; By PO-coherence
(Definition~\ref{def:po-coherent}), there exists a type $\sigma$ with
$C \in \sigma$ and $\PO \in \bigcap_j \Safe(\tau_j, \sigma)$.
The model guarantees that $\sigma$ is a realized type (it appears as
the type of a $\PO$-witness at some chain position).  Create a
quotient witness of type $\sigma$ with $\rho(k_\alpha, w') = \PO$.

\medskip
The total number of designated witnesses is at most
$K \cdot |\cl(C_0)|$ (one per demand per kernel), so
$L \le D \cdot |\cl(C_0)|$.

\medskip\noindent\textit{Step 4: Witness-to-other relations and
V5.}\;
For each designated witness $w'$ and each other quotient node $e$:
the model provides a globally consistent relation $\rho(w, e)$ for
the original witness element.  By the AC-survival argument (below),
these values survive arc-consistency enforcement.

The AC-survival argument: each off-chain witness $w$ in the model
has actual relations to all other elements.  For any existing node
$e$, $\rho(w, e) \in D(w, e)$ (the initial domain from
\eqref{eq:v6-domain}), because $\rho(w, e) \in
\comp(\rho(w, \text{parent}), \rho(\text{parent}, e))$ (composition)
and $\rho(w, e) \in \Safe(\tp(w), \tp(e))$ (model type-safety).
Arc-consistency cannot remove $\rho(w, e)$ because every mediator
$e'$ satisfies $\rho(w, e') \in D(w, e')$ and
$\rho(w, e') \in \comp(\rho(w, e), \rho(e, e'))$ (global
consistency).  So all domains remain non-empty.

\medskip\noindent\textit{Step 5: $V_{\mathrm{safe}}$ for
kernel-to-regular edges.}\;
Each regular node inherits its model relations.  By
Lemma~\ref{lem:phase-uniform}, the stabilized kernel-to-regular
relation satisfies $V_{\mathrm{safe}}$.

\medskip\noindent\textit{Step 6: Off-chain $\PP$-witnesses (V6).}\;
For each $\exists\PP.C \in \Union(\Desc_\alpha)$ not satisfied
within $\Desc_\alpha$: the model provides a $\PP$-witness $w$
at some chain position $i_n > M$.  By
Lemma~\ref{lem:backward-safe} (with $R = \PP$),
$\PP \in \Safe(\tau_j, \tp(w))$ for every phase $\tau_j$.
The AC-survival argument ensures the V6 path-consistency check
passes.

\medskip\noindent\textit{Size bound.}\;
$K + L \le D + D \cdot |\cl(C_0)| = D\,(1 + |\cl(C_0)|)$,
where $D \le 2^{2^{|\cl(C_0)|}}$.  This is doubly exponential but
\textbf{computable}, with no circular dependency.
\end{proof}


%% ============================================================
\section{Soundness: quotient implies model}\label{sec:soundness}

\begin{lemma}[Chain-unfolding lift]\label{lem:lift}
Let $\mathcal{N}$ be a CC atomic RCC5 network on kernel and regular
nodes.  For each kernel $k_\alpha$ with descriptor
$\Desc_\alpha = (\tau_1, \ldots, \tau_p)$, replace $k_\alpha$ by an
infinite chain $d_1^{(\alpha)}, d_2^{(\alpha)}, \ldots$ with:
\begin{itemize}
\item $\rho(d_i^{(\alpha)}, d_j^{(\alpha)}) = \PP$ for $i < j$,
  $\EQ$ for $i = j$;
\item $\rho(d_i^{(\alpha)}, e) = \rho(k_\alpha, e)$ for every
  non-chain node $e$ and all $i$;
\item $\rho(d_i^{(\alpha)}, d_j^{(\beta)}) =
  \rho(k_\alpha, k_\beta)$ for distinct chains $\alpha \ne \beta$.
\end{itemize}
Then the unfolded network $\mathcal{N}'$ is CC.
\end{lemma}

\begin{proof}
Every triple in $\mathcal{N}'$ falls into one of six cases:

\emph{Case~1: within-chain $(d_i, d_j, d_k)$, $i < j < k$.}\;
$\PP \in \comp(\PP, \PP) = \{\PP\}$. \checkmark

\emph{Case~2: two same-chain elements $d_i, d_j$ ($i < j$) and
external $e$.}\;
$\rho(d_i, e) = \rho(d_j, e) = R$ (constant interface).
Need: $R \in \comp(\PPI, R)$.  Holds by
Theorem~\ref{thm:self-absorbing}(c). \checkmark

\emph{Case~3: one chain element $d_i^{(\alpha)}$ and two externals
$e_1, e_2$.}\;
Reduces to triple $(k_\alpha, e_1, e_2)$, CC in $\mathcal{N}$.
\checkmark

\emph{Case~4: elements from two chains and a third node.}\;
Reduces to $(k_\alpha, k_\beta, e)$, CC in $\mathcal{N}$. \checkmark

\emph{Case~5: three elements from three chains.}\;
Reduces to $(k_\alpha, k_\beta, k_\gamma)$, CC in $\mathcal{N}$.
\checkmark

\emph{Case~6: two same-chain elements and one from another chain.}\;
Reduces to Case~2 with $e = d_j^{(\beta)}$. \checkmark
\end{proof}

\begin{corollary}[V6 lifts to the unfolded model]\label{cor:lift}
If the finite V6 extension network is path-consistent, then the
unfolded network is also path-consistent.\qed
\end{corollary}

\begin{theorem}[Quotient realization]\label{thm:soundness}
If a valid two-tier quotient $\mathcal{Q}$ exists for $C_0$, then
$C_0$ is satisfiable.
\end{theorem}

\begin{proof}
We construct a model $\interp' \models C_0$.

\medskip\noindent\textit{Step 1: Unfold descriptors into chains.}\;
For each $\Desc_\alpha = (\tau_1, \ldots, \tau_p)$, create an
infinite chain $d_1^{(\alpha)}, d_2^{(\alpha)}, \ldots$ with
$\tp(d_i^{(\alpha)}) = \tau_{((i-1) \bmod p) + 1}$ and
$\rho(d_i, d_j) = \PP$ for $i < j$, $\EQ$ for $i = j$.

\medskip\noindent\textit{Step 2: Assign constant kernel
interfaces.}\;
$\rho(d_i^{(\alpha)}, e) = \rho(k_\alpha, e)$ for every non-chain
node $e$.

\medskip\noindent\textit{Step 3: CC of the unfolded base.}\;
By Lemma~\ref{lem:lift}.

\medskip\noindent\textit{Step 4: Add off-chain witnesses.}\;
V6 ensures path-consistency.  By Corollary~\ref{cor:lift}, this
lifts to the unfolded network.  By
Theorem~\ref{thm:full-tractability}, a consistent atomic refinement
exists.  Every value in the refined domains is in $\Safe$, so the
chosen atomic relations respect all universal constraints.

\medskip\noindent\textit{Step 5: Demand satisfaction.}\;
\begin{itemize}
\item \textbf{$\exists\PP$ within chains:} If $C \in \tau_j$ for
  some $j$, the next occurrence of $\tau_j$ on the chain witnesses
  the demand.  Otherwise, V6 provides an off-chain $\PP$-witness.
\item \textbf{Non-$\PP$ demands of chain elements:} For
  $\exists R.C \in \tau_i$ with $R \ne \PP$: the designated witness
  $w$ (T4) has $\rho(k_\alpha, w) = R$, so
  $\rho(d_j^{(\alpha)}, w) = R$ (Step~2), witnessing the demand.
\item \textbf{Universal constraints:} For $\forall R.C \in \tau_i$
  and an external $e$ with $\rho(d_j, e) = R$:
  $V_{\mathrm{safe}}$ guarantees
  $R \in \Safe(\tau_i, \tp(e))$, so $C \in \tp(e)$.
\end{itemize}

\medskip\noindent\textit{Step 6: Global consistency.}\;
The full atomic network is CC (Steps~3--4).  By the patchwork
property, this CC atomic network is globally consistent and
realizable, giving $\interp' \models C_0$.
\end{proof}


%% ============================================================
\section{Main result}\label{sec:decidability}

\begin{theorem}[Decidability of $\ALCIRCC{5}$, PO-coherent fragment]
\label{thm:decidability}
Concept satisfiability in $\ALCIRCC{5}$ is decidable for PO-coherent
concepts.
\end{theorem}

\begin{proof}
Given $C_0$:
\begin{enumerate}
\item Compute $f(|C_0|) = D\,(1 + |\cl(C_0)|)$ with
  $D = 2^{2^{|\cl(C_0)|}}$.
\item Enumerate all two-tier quotients of size $\le f(|C_0|)$.
\item For each, check validity: (V1)--(V4) for descriptors,
  PO-coherence, (V5) CC, ($V_{\mathrm{safe}}$) phasewise safety,
  (V6) off-chain path-consistency, (T4)--(T5) designated witnesses.
\item Accept iff any valid quotient is found.
\end{enumerate}
Correctness: no false negatives
(Theorem~\ref{thm:completeness}), no false positives
(Theorem~\ref{thm:soundness}).  All steps are computable.
\end{proof}


%% ============================================================
\section{The $\PO$ gap}\label{sec:po-gap}

\subsection{Why $\PO$ is different}

For $R \in \{\DR, \PP\}$, backward forcing
(Lemma~\ref{lem:backward}) guarantees that a witness at a chain
position past the exhaustion index $M$ is $R$-related to
\emph{all} earlier chain positions.  This yields all-phase
type-safety (Lemma~\ref{lem:backward-safe}), enabling a
constant-$R$ kernel interface.

For $R = \PPI$, forward absorption ($\comp(\PPI, \PPI) = \{\PPI\}$)
ensures that the stabilized value of any $\PPI$-witness \emph{is}
$\PPI$, and Lemma~\ref{lem:phase-uniform} gives all-phase safety.

For $R = \PO$, \emph{neither} property holds
(Remark~\ref{rem:no-backward-po}).  The composition table permits
$\PO$ at one chain position and a different relation ($\DR$, $\PP$,
or $\PPI$) at another.

\subsection{A realizable PO-incoherent descriptor}

\begin{proposition}\label{prop:po-incoherent}
There exists a satisfiable $\ALCIRCC{5}$ concept $C_0$ whose
model-extracted period descriptor $\Desc = (\tau_A, \tau_B)$
satisfies $\exists\PO.A \in \tau_A$ and
$\forall\PO.\neg A \in \tau_B$, yet no element in the model has
stabilized relation $\PO$ to the chain.
\end{proposition}

\begin{proof}[Construction]
Let the chain be $d_0(\tau_A)\;\PP\;d_1(\tau_B)\;\PP\;
d_2(\tau_A)\;\PP\;d_3(\tau_B)\;\PP\;\cdots$.  For each $k \ge 0$,
create a witness $w_k$ with $A \in \tp(w_k)$ and:
\[
\rho(d_i, w_k) = \begin{cases}
\DR & \text{if } i \le 2k - 1, \\
\PO & \text{if } i = 2k, \\
\PPI & \text{if } i \ge 2k + 1.
\end{cases}
\]
The transitions $\DR \to \PO \to \PPI$ are valid by
Lemma~\ref{lem:forward}.  Backward forcing gives $\DR$ at all
positions before the $\PO$ event.

\emph{$\forall\PO.\neg A$ at $\tau_B$-positions:}\; At
$d_{2k+1}$ (type $\tau_B$), the relation to $w_k$ is $\PPI$, to
$w_l$ for $l < k$ is $\PPI$, and to $w_l$ for $l > k$ is $\DR$.
No element is $\PO$-related to $d_{2k+1}$, so
$\forall\PO.\neg A$ is vacuously satisfied.

\emph{Mutual witness consistency:}\; For $j < k$:
$\rho(w_j, w_k) \in \comp(\inv(\rho(d_{2j+1}, w_j)),
\rho(d_{2j+1}, w_k)) = \comp(\PP, \DR) = \{\DR\}$.  So
$\rho(w_j, w_k) = \DR$ for all pairs, which is
composition-consistent.

Every $w_k$ stabilizes to $\PPI$.  The concept is satisfiable, but
no element has stabilized $\PO$.
\end{proof}

\begin{corollary}
The constant-interface quotient architecture is incomplete for
PO-incoherent descriptors: a constant-$\PO$ witness at the kernel
would be $\PO$-related to $\tau_B$-positions, violating
$\forall\PO.\neg A$.
\end{corollary}

\begin{remark}[The $\DR$ analogue is impossible]
The analogous descriptor with $\exists\DR.A \in \tau_A$ and
$\forall\DR.\neg A \in \tau_B$ is \emph{unrealizable}: any
$\DR$-witness at a $\tau_A$-position is $\DR$ at all earlier
$\tau_B$-positions (backward forcing), violating
$\forall\DR.\neg A$.  This is why the $\DR$ case has no gap.
\end{remark}

\subsection{The scope of the gap}

The $\PO$ gap affects \emph{only} descriptors that are
PO-incoherent in the sense of Definition~\ref{def:po-coherent}.
Specifically, the gap arises only when:
\begin{enumerate}
\item Some phase $\tau_j$ contains $\exists\PO.C$.
\item Some other phase $\tau_k$ contains a $\forall\PO$-universal
  that is incompatible with \emph{every} witness type satisfying $C$.
\end{enumerate}
This is a relatively specific configuration.  For the vast majority
of $\ALCIRCC{5}$ concepts---including all those without
$\forall\PO$ restrictions, those where $\forall\PO$ universals are
uniform across phases, and those where the witness type is compatible
with all phases---the proof is complete.


%% ============================================================
\section{Complexity}\label{sec:complexity}

The quotient size bound is $D\,(1 + |\cl(C_0)|)$ with
$D \le 2^{2^{|\cl(C_0)|}}$.  This is doubly exponential; we do not
claim \EXPTIME{} membership.  An \EXPTIME{} bound (matching the
$\ALCI$ lower bound) would require showing the quotient size is
singly exponential, which we leave open.


%% ============================================================
\section{Discussion}\label{sec:discussion}

\subsection{The role of full RCC5 tractability}

The soundness proof relies on Theorem~\ref{thm:full-tractability}:
path-consistency decides satisfiability for the entire RCC5 algebra.
Without it, we would need to verify global consistency directly.
Full tractability collapses this gap: local consistency (checkable)
implies global consistency (needed).

\subsection{Why constant interfaces work for $\DR$, $\PP$, $\PPI$}

The backward forcing property ($\comp(\PP, \DR) = \{\DR\}$ and
$\comp(\PP, \PP) = \{\PP\}$) is the key algebraic fact.  When a
$\DR$-witness (resp.\ $\PP$-witness) exists at chain position $i_n$,
backward forcing makes it $\DR$ (resp.\ $\PP$) at \emph{all}
earlier positions.  Since all recurrent types appear at earlier
positions, all-phase type-safety is automatic
(Lemma~\ref{lem:backward-safe}).

For $\PPI$, forward absorption ($\comp(\PPI, \PPI) = \{\PPI\}$)
plays the dual role: the stabilized value is always $\PPI$, and
Lemma~\ref{lem:phase-uniform} gives all-phase safety.

\subsection{The algebraic root of the $\PO$ gap}

$\PO$ is the unique base relation with \emph{neither} backward
forcing nor forward absorption.  The composition entries
$\comp(\PP, \PO) = \{\DR, \PO, \PP\}$ (not $\{\PO\}$) and
$\comp(\PPI, \PO) = \{\PO, \PPI\}$ (not $\{\PO\}$) are the
algebraic root.  Furthermore, the cycle consistency condition
$\PO \in \comp(\PPI, \PPI) = \{\PPI\}$ fails, so even
phase-indexed interfaces cannot make $\PO$ recur periodically
along a $\PP$-chain once interrupted by $\PPI$.

\subsection{Potential routes to closing the $\PO$ gap}

\begin{enumerate}
\item \textbf{Proof that PO-incoherent descriptors are unrealizable.}\;
  If one could show that $\exists\PO.A \in \tau_j$ and
  $\forall\PO.\neg A \in \tau_k$ cannot coexist in any model's
  $T_\infty$, the gap would vanish.  Proposition~\ref{prop:po-incoherent}
  refutes this unless additional semantic constraints are identified.
\item \textbf{A non-constant quotient architecture.}\; Replacing
  constant kernel interfaces with per-position interfaces would
  handle $\PO$ but would make the quotient infinite (or require a
  different finiteness argument).
\item \textbf{Hybrid approach.}\; Use the two-tier quotient for
  $\DR/\PP/\PPI$ demands and a different mechanism (e.g., type
  elimination~\cite{Claude2026}) for $\PO$ demands.
\end{enumerate}

\subsection{Extension to $\ALCIRCC{8}$}

For RCC8, $\PP$ splits into $\mathsf{TPP}$ and $\mathsf{NTPP}$.
$\mathsf{NTPP}$ is transitive, but RCC8 is \emph{not} fully
tractable.  The decidability of $\ALCIRCC{8}$ remains open.


%% ============================================================
\section{Conclusion}\label{sec:conclusion}

We have proved that concept satisfiability in $\ALCIRCC{5}$ is
decidable for the PO-coherent fragment.  The proof rests on four
pillars:
\begin{enumerate}
\item \textbf{Algebraic:} Backward forcing of $\DR$ and $\PP$;
  forward absorption of $\PPI$; self-absorption of $\PP$.
\item \textbf{Semantic:} The phase-uniform type-safety lemma
  (Lemma~\ref{lem:phase-uniform}) and the backward-forcing safety
  lemma (Lemma~\ref{lem:backward-safe}).
\item \textbf{Constructive:} One kernel per descriptor, yielding
  a computable size bound $D\,(1 + |\cl(C_0)|)$ with no circularity.
\item \textbf{Constraint-theoretic:} Full RCC5 tractability bridges
  local consistency (V6) to global consistency (the model).
\end{enumerate}

The chain-unfolding lift lemma (Lemma~\ref{lem:lift}) ensures the
finite V6 certificate governs the infinite unfolded model.

The decidability of $\ALCIRCC{5}$ in full generality remains open
due to the $\PO$ gap (Proposition~\ref{prop:po-incoherent}).  The
$\PO$ obstacle is precisely localized: it arises only for
PO-incoherent descriptors where different phases carry conflicting
$\forall\PO$ universals.  Whether this obstacle can be overcome
within the two-tier framework or requires a fundamentally different
approach is the key remaining question.


\begin{thebibliography}{99}

\bibitem{Wessel2003}
Michael Wessel.
\newblock \emph{Qualitative Spatial Reasoning with the ALCIRCC Family
  -- First Results and Unanswered Questions}.
\newblock Technical Report FBI-HH-M-324/03, University of Hamburg,
  2002/2003.

\bibitem{Claude2026}
Claude (Anthropic).
\newblock \emph{On the Decidability of ALCIRCC5 and ALCIRCC8:
  Quasimodels Meet the Patchwork Property}.
\newblock Discussion draft, March 2026.

\bibitem{ClaudeResponse2026}
Claude (Anthropic).
\newblock \emph{Response to GPT-5.4's Review of the Two-Tier Quotient
  Paper}.
\newblock Discussion note, April 2026.

\bibitem{GPT2026}
GPT-5.4 Pro.
\newblock \emph{A Contextual Tableau Calculus for
  $\mathcal{ALCI}_{\mathit{RCC5}}$}.
\newblock Discussion draft, 2026.

\bibitem{GPTReview2026}
GPT-5.4 Pro.
\newblock \emph{Technical Response to ``Decidability of
  $\mathcal{ALCI}_{\mathit{RCC5}}$ via Two-Tier Quotient
  Construction''}.
\newblock First review, April 2026.

\bibitem{GPTReview2026b}
GPT-5.4 Pro.
\newblock \emph{Second Technical Response to the Revised Manuscript}.
\newblock Second review, April 2026.

\bibitem{GPTReview2026c}
GPT-5.4 Pro.
\newblock \emph{Technical Response to the Third Revision}.
\newblock Third review, April 2026.

\bibitem{ClaudeFW2026}
Claude (Anthropic).
\newblock \emph{On the Failure of Finite-Width Extraction for
  $\mathcal{ALCI}_{\mathit{RCC5}}$}.
\newblock Discussion note, March 2026.

\bibitem{ClaudeTriangle2026}
Claude (Anthropic).
\newblock \emph{Triangle Types and the Extension Gap in
  $\mathcal{ALCI}_{\mathit{RCC5}}$}.
\newblock Discussion draft, March 2026.

\bibitem{RenzNebel1999}
Jochen Renz and Bernhard Nebel.
\newblock On the complexity of qualitative spatial reasoning:
  A~maximal tractable fragment of the Region Connection Calculus.
\newblock \emph{Artificial Intelligence}, 108(1--2):69--123, 1999.

\bibitem{Renz1999}
Jochen Renz.
\newblock Maximal tractable fragments of the Region Connection
  Calculus: A~complete analysis.
\newblock In \emph{Proc.\ IJCAI}, pages 448--454, 1999.

\bibitem{RandellCuiCohn1992}
David~A. Randell, Zhan Cui, and Anthony~G. Cohn.
\newblock A spatial logic based on regions and connection.
\newblock In \emph{Proc.\ KR}, pages 165--176, 1992.

\end{thebibliography}

\end{document}
